\chapter{Shannon and von Neumann}
\label{ch:shannon}

Two people stand behind every page of this book. John von Neumann\index{von Neumann, John} described the architecture of the stored program computer, invented Monte Carlo with Ulam\index{Ulam, Stanislaw}, gave the rejection method\index{rejection method}, showed how to extract fair coins from biased ones, and worried in print about whether arithmetic could produce randomness at all. Claude Shannon\index{Shannon, Claude} founded information theory, proved by a random construction that good codes exist, counted Boolean circuits to show that most functions are hard, and wrote the first paper on programming a computer to play chess. This chapter records what they said, because the Monte Carlo program of Part IV is, in a direct line, theirs.

\section{Von Neumann and the architecture of computation}

Von Neumann's First Draft of a Report on the EDVAC of 1945 set out the architecture in which program and data share a memory, so that a machine can modify its own instructions, which is Turing\index{Turing, Alan}'s universal machine made of vacuum tubes. The report was circulated under his name alone, to the lasting resentment of Eckert and Mauchly, and it shaped the design of nearly every computer since. For complexity theory its importance is that it made the universal machine real and therefore made the question of running time an engineering question as well as a mathematical one.

In 1956 von Neumann published Probabilistic Logics and the Synthesis of Reliable Organisms from Unreliable Components~\cite{vonNeumann56}, which asked how a reliable computation can be built from components that fail with probability $\epsilon$. His answer, redundancy with majority voting, is the ancestor of error correcting computation and of the fault tolerance thresholds of quantum computing. The paper also contains an early statement of a theme of this book: that a probabilistic process can compute a deterministic answer with high confidence, and that the confidence can be amplified by repetition at logarithmic cost. The amplification lemma of Chapter~\ref{ch:pseudo} is anticipated in it.

\section{Von Neumann and Monte Carlo}

The Monte Carlo method was conceived by Stanislaw Ulam in 1946, while convalescing and playing solitaire, as a way to estimate probabilities by repeated trial, and developed with von Neumann for neutron diffusion calculations at Los Alamos. Von Neumann's letter to Richtmyer of March 1947 is the founding document: it describes how to sample the paths of neutrons through fissile material by generating random numbers for each collision, and proposes running the calculation on the ENIAC. The name was Metropolis\index{Metropolis algorithm}'s. The method was the first use of a computer to simulate a physical process by random sampling, and it is the direct ancestor of every algorithm in this book.

Von Neumann made three technical contributions to Monte Carlo that we use. The first is the rejection method, in his 1951 paper Various Techniques Used in Connection with Random Digits~\cite{vonNeumann51}: to sample from a density $f$, draw a candidate from an easier density $g$ with $f \le c g$, and accept it with probability $f / (c g)$. The expected number of candidates per accepted sample is $c$. This is the simplest level set sampler\index{level set sampler}, and it is exactly the naive method for the queens, with $g$ uniform on permutations, $f$ uniform on solutions, and $c = 1/\rho_N = e^{0.942 N}$; it is correct and exponentially slow, and the whole of Part III is about doing better.

The second is the extraction of fair bits from a biased coin, in the same paper. Flip the coin twice; if the results differ, output the first; if they agree, discard and repeat. The output is an unbiased bit regardless of the bias, at an expected cost of $1 / (p(1-p))$ flips. This is the first randomness extractor, and the theory of extractors, developed by Nisan, Zuckerman, Trevisan and others in the 1990s, is the bridge between the pseudorandom generator\index{pseudorandom generator}s of Chapter~\ref{ch:pseudo} and the physical sources of randomness that a real computer uses.

The third is the remark, in the same paper, that has become the most quoted sentence in the subject: anyone who considers arithmetical methods of producing random digits is, of course, in a state of sin. Von Neumann proposed the middle square method anyway, squaring a number and taking the middle digits, and knew it was poor. His point was that a deterministic procedure cannot produce randomness in the full sense, and that what it produces must be judged by whether it serves the purpose, which is the pragmatic definition of pseudorandomness that Yao and Blum and Micali made rigorous thirty years later: a sequence is random enough if no efficient test can tell it from random. The hardness versus randomness programme is the theory of the state of sin, and the question $\Pc = \BPP$ is the question whether the sin has consequences.

\section{Von Neumann and the inversion of positive operators}

One further strand connects von Neumann to the vertical likelihood method. In their work on the ENIAC calculations, von Neumann and Ulam proposed, in a scheme published and credited to them by Forsythe and Leibler in 1950, estimating the entries of the inverse of a matrix, or more generally the solution of a positive linear inverse problem, by sampling random walks whose transition probabilities are read off from the matrix and whose weighted path sums are unbiased for the required quantities. Vardi and Lee (1993) treated positive linear inverse problems by a maximum likelihood iteration of EM type, and one of us, in a 1993 discussion with Robert~\cite{PolsonRobert}, placed it in the Bayesian framework and showed that the Monte Carlo inversion of a positive operator is a form of importance sampling whose variance is controlled by the spectral properties of the operator. The level set samplers of Part III are random walks of the same kind on the configuration space, and their mixing is governed by the spectral gap\index{spectral gap} of the same kind of operator.

\section{Shannon and information}

Shannon's A Mathematical Theory of Communication of 1948~\cite{Shannon48} defined the entropy of a source, the capacity of a channel, and proved that reliable communication is possible at any rate below capacity and impossible above it. The proof of the direct part is an early great application of the probabilistic method: choose a code at random, show that the expected error probability is small, conclude that a good code exists. Shannon did not exhibit a good code, and explicit asymptotically good codes appeared only with Justesen in 1972; the random construction proved existence without construction, exactly as Adleman\index{Adleman's theorem}'s argument in Chapter~\ref{ch:pseudo} proves that a good random string exists without finding it. The theory of explicit codes, from Reed and Solomon through Justesen to the polar and LDPC codes of the present, is the derandomization of Shannon's argument, and it took decades, with explicit capacity achieving polar codes dating from 2009.

Entropy is the quantity that governs the counting in this book. Simkin\index{Simkin's theorem}'s upper bound in Chapter~\ref{ch:simkin} is a chain rule for entropy; the density of states\index{density of states} is $e^{\text{entropy}}$ at each energy; the ladder length is the entropy difference between the uniform distribution and the uniform distribution on solutions; and the rate of a rare event is, by Sanov, a relative entropy. Shannon's source coding theorem, that entropy is the compression limit, has a direct reading: a uniform solution of the queens problem can be described in $\log_2 Q(N) \approx N \log_2 N - 2.8 N$ bits, against $\log_2 N! \approx N \log_2 N - 1.44 N$ for a uniform permutation, and the difference of $1.36 N$ bits is the number of bits, paid in rejected samples, that the vertical ladder must spend to get from one to the other.

\section{Shannon and circuits}

Shannon's master's thesis of 1937 showed that Boolean algebra describes switching circuits, which is why we call them Boolean circuits. His paper of 1949 on the synthesis of two terminal switching circuits~\cite{Shannon49} contains the counting argument of Chapter~\ref{ch:circuits}: almost all Boolean functions of $n$ variables require about $2^n / n$ switches. It is the first circuit lower bound, it is nonconstructive, and it has not been matched for any explicit function in the more than seventy five years since. Shannon's argument proves that hard functions exist; the inability to point to an explicit one is the reason $\Pc$ versus $\NP$ is open.

\section{Shannon and secrecy}

Shannon's Communication Theory of Secrecy Systems, written in 1945 and published in 1949~\cite{Shannon49b}, proved that the one time pad is perfectly secure and that perfect security requires a key as long as the message. Every practical cipher is therefore imperfect, and its security must rest on the computational difficulty of recovering the key, which is a question about $\NP$: the key is a short certificate that the ciphertext decrypts to a given plaintext. Shannon introduced the unicity distance, the length of ciphertext beyond which the key is in principle determined, and distinguished in principle recoverability from practical recoverability, which is the distinction between information theoretic and computational security and the reason that cryptography lives in Minicrypt and Cryptomania rather than in information theory. Diffie and Hellman's 1976 paper, which founded public key cryptography, cites Shannon and makes the computational assumption explicit.

\section{Shannon and chess}

Shannon's Programming a Computer for Playing Chess of 1950~\cite{Shannon50} is the first paper on computer chess. He estimated the number of possible games at about $10^{120}$, the Shannon number, by assuming about $30$ legal moves per position and $40$ moves per side, and concluded that exhaustive search was out of the question. He proposed two strategies: type A, searching all moves to a fixed depth and evaluating the leaves with a static function, and type B, searching selectively along plausible lines as a human does. Every chess engine since has been a mixture, with alpha beta pruning making type A far stronger than Shannon expected and the neural evaluations of AlphaZero reviving type B.

The queens problem is not chess, but it is chess pieces, and Shannon's estimate is the right comparison for the counting problem. There are $N^N$ placements and $N!$ permutations and $Q(N)$ solutions, and for $N = 27$ the exact count of $2.3 \times 10^{17}$ solutions among $1.1 \times 10^{28}$ permutations took a year of hardware. For $N = 100$ there are about $10^{116}$ solutions among $10^{158}$ permutations, and no exhaustive method will ever count them. Shannon's conclusion for chess, that one must search selectively, is the conclusion for the queens, and the vertical ladder is a type B strategy: it searches the staircase from the top, following the plausible lines, which are the level sets, and prunes the rest.

\section{What they would say now}

Von Neumann died in 1957 and Shannon in 2001, and neither contributed to the theory of $\NP$ completeness, though Shannon lived through it. We may guess at their reactions. Von Neumann, who proved the minimax theorem, the Monte Carlo method and the rejection sampler, would have recognised the vertical ladder as a rejection sampler run in stages, and would have asked at once for the mixing time\index{mixing time} of the chain at the bottom. Shannon, who proved existence by random coding and spent the rest of his career watching others try to make it explicit, would have recognised the derandomization problem as the problem he set in 1948, and would not have been surprised that it is still open. Both would have agreed that the question of whether a coin helps computation is a question about physics as much as mathematics, and both would have been pleased, we think, that it can be stated on a chessboard.
